% !TEX program = pdflatex
% Relatório executivo — Comparativo Liebherr × Hyundai
\documentclass[11pt,a4paper]{article}

\usepackage[utf8]{inputenc}
\usepackage[T1]{fontenc}
\usepackage[brazil]{babel}
\usepackage{lmodern}
\usepackage{geometry}
\geometry{margin=2.5cm}

\usepackage{graphicx}
\usepackage{float}
\usepackage{booktabs}
\usepackage{array}
\usepackage{tabularx}
\usepackage{longtable}
\usepackage{pdflscape}
\usepackage{rotating}
\usepackage{hyperref}
\usepackage{xcolor}
\usepackage{enumitem}

\hypersetup{
  colorlinks=true,
  linkcolor=black,
  urlcolor=blue,
  pdftitle={Comparativo Liebherr x Hyundai},
}

% Figuras copiadas para ./figuras/ (sincronizar_figuras.py ou Passo 9 do notebook)
\graphicspath{{figuras/}}
\newcommand{\figdir}{figuras}

% Preenchidos por dados.tex (exportados do notebook)
\newcommand{\reportdate}{08/06/2026}
\newcommand{\daterange}{2025/01 a 2026/03}
\newcommand{\reportauthor}{Julio}
\newcommand{\reportarea}{Qualidade - Pilares}
\newcommand{\reportcredit}{Análises feitas por Julio; Qualidade - Pilares}



\begin{document}

\begin{titlepage}
  \centering
  \vspace*{2cm}

  {\LARGE\bfseries Comparativo de Custos de Máquinas\par}
  \vspace{0.6cm}
  {\Large Liebherr \times\ Hyundai\par}
  \vspace{1.2cm}

  {\large Relatório executivo — Controladoria\par}
  \vspace{0.4cm}
  {\normalsize Período: \daterange\par}
  \vspace{0.2cm}
  {\normalsize Parque: 7 Liebherr · 16 Hyundai\par}
  \vspace{0.3cm}
  {\normalsize Emitido em: \reportdate\par}

  \vfill

  {\small
    Documento preparado para apoio à gestão financeira.\\
    Comparativos entre marcas em custo médio por máquina do parque.\\
    Não inclui detalhamento técnico da extração de dados.
  \par}

  \vspace{0.8cm}
  {\normalsize\textit{\reportcredit}\par}

  \vspace{1.5cm}
\end{titlepage}

\tableofcontents
\newpage


\section*{Resumo executivo}
\addcontentsline{toc}{section}{Resumo executivo}

Comparativo de custos das frotas \textbf{Liebherr} (7 máquinas) e \textbf{Hyundai} (16 máquinas)
no período \daterange. Gráficos Liebherr $\times$ Hyundai usam o \textbf{custo médio por máquina}
(gasto total do parque $\div$ 7 ou $\div$ 16), não totais brutos --- assim a comparação não favorece
a marca com mais equipamentos. Ranking por identificador único (\texttt{maq\_id}).

\begin{itemize}[leftmargin=*]
  \setlength{\itemsep}{0.35em}
  \item \textbf{Custo médio/máquina (comparativo justo):} LIEBHERR custa \textbf{42.8\%} a mais (R\$ 124,072/máq vs R\$ 86,871/máq; parque: 7 vs 16 máquinas).
  \item \textbf{Volume absoluto (contexto):} HYUNDAI tem maior gasto total (R\$ 1,399,720 vs R\$ 870,203) — reflexo do tamanho do parque.
  \item \textbf{2025 — HYUNDAI:} volume R\$ 1,220,373 (custo médio R\$ 75,649/máquina).
  \item \textbf{2025 — LIEBHERR:} volume R\$ 733,840 (custo médio R\$ 104,834/máquina).
  \item \textbf{2026 — HYUNDAI:} volume R\$ 179,347 (custo médio R\$ 11,209/máquina).
  \item \textbf{2026 — LIEBHERR:} volume R\$ 136,362 (custo médio R\$ 19,237/máquina).
  \item \textbf{Tipo dominante (R\$/máq):} LIEBHERR $\rightarrow$ 7.1.2 MANUTENÇÃO DE VEÍCULOS/MAQUINAS (R\$ 52,641/máq); HYUNDAI $\rightarrow$ 7.1.2 MANUTENÇÃO DE VEÍCULOS/MAQUINAS (R\$ 46,949/máq).
  \item \textbf{Mix manutenção (7.1.2):} LIEBHERR 42.4\% vs HYUNDAI 54.1\% do volume.
  \item \textbf{Mix combustível (7.1.22+7.1.4):} LIEBHERR 25.2\% vs HYUNDAI 35.6\%.
  \item \textbf{Mix peças (7.1.1):} LIEBHERR 30.5\% vs HYUNDAI 8.2\%.
  \item \textbf{Concentração LIEBHERR:} top 3 máquinas = 82\% do gasto (líder: EHL0043).
  \item \textbf{Concentração HYUNDAI:} top 3 máquinas = 38\% do gasto (líder: EHH0044).
  \item \textbf{Unidade com maior custo:} Maringá — R\$ 556,716 (5 máq.; R\$ 111,343/máq).
  \item \textbf{Top 3 unidades:} Maringá (R\$ 556,716); Dourados (R\$ 429,208); Prudente (R\$ 349,489) — 59\% do gasto do parque.
  \item \textbf{Maior intensidade (R\$/máq):} Dourados — R\$ 143,069/máq (3 máq.).
  \item \textbf{Na unidade líder (Maringá):} LIEBHERR R\$ 348,601 vs HYUNDAI R\$ 208,115 (predomina LIEBHERR).
\end{itemize}


\section{Indicadores principais}
\addcontentsline{toc}{section}{Indicadores principais}

\begin{table}[H]
\centering
\small
\begin{tabular}{lrr}
\toprule
Indicador & HYUNDAI & LIEBHERR \\
\midrule
Lançamentos (marca) & 1.077 & 386 \\
Lançamentos (parque) & 1.063 & 384 \\
Máquinas no parque & 16 & 7 \\
Máquinas com lançamento & 16 & 7 \\
Resultado líquido (valor conta) & R\$\,545.063 & R\$\,668.859 \\
Volume de gasto (|valor|) & R\$\,1.399.720 & R\$\,870.203 \\
Volume parque (|valor|) & R\$\,1.389.933 & R\$\,868.503 \\
Despesas & R\$\,972.392 & R\$\,769.531 \\
Ajustes / créditos & R\$\,427.329 & R\$\,100.672 \\
Ticket médio & R\$\,1.300 & R\$\,2.254 \\
Custo médio/máquina (parque) & R\$\,86.871 & R\$\,124.072 \\
\bottomrule
\end{tabular}
\caption{Indicadores comparativos por marca.}
\end{table}


\section{Custo médio por máquina (por ano)}
\addcontentsline{toc}{section}{Custo médio por máquina (por ano)}

Não se soma 2025 e 2026 num único total: cada ano é analisado à parte. O valor é a
\textbf{média simples}: gasto total do parque dividido pela quantidade oficial de máquinas (7 ou 16).

Os gráficos deste relatório \textbf{excluem} a conta \texttt{7.4.5} (ISS) --- imposto municipal, não custo operacional da máquina. Tabelas de KPI e cadastro mantêm a base completa.

\begin{figure}[H]
  \centering
  \includegraphics[width=0.88\linewidth]{\figdir/01_volume_por_ano.png}
  \caption{Custo médio de cada máquina do parque em cada ano (volume total $\div$ 7 ou 16).}
\end{figure}

\section{Composição por tipo de despesa}

\begin{figure}[H]
  \centering
  \includegraphics[width=\linewidth]{\figdir/02_mix_contas_pct.png}
  \caption{Mix percentual por tipo de conta (top rubricas) --- percentual \textit{dentro} de cada marca.}
\end{figure}

\begin{figure}[H]
  \centering
  \includegraphics[width=\linewidth]{\figdir/03_volume_por_conta.png}
  \caption{Gasto por conta em custo médio por máquina do parque.}
\end{figure}

\section{Evolução mensal}

\begin{figure}[H]
  \centering
  \includegraphics[width=\linewidth]{\figdir/04_serie_mensal_2025.png}
  \caption{Custo médio mensal por máquina do parque --- 2025.}
\end{figure}

\begin{figure}[H]
  \centering
  \includegraphics[width=\linewidth]{\figdir/04_serie_mensal_2026.png}
  \caption{Custo médio mensal por máquina do parque --- 2026.}
\end{figure}

\section{Localização do parque}
\addcontentsline{toc}{section}{Localização do parque}

Distribuição informada pelo diretor financeiro, detalhada por identificador de máquina (\texttt{maq\_id}) e referência de centro de custo na base.

\begin{table}[H]
\centering
\small
\begin{tabular}{lrrr}
\toprule
Local & LIEBHERR & HYUNDAI & Total \\
\midrule
Prudente & 1 & 2 & 3 \\
Maringá & 2 & 3 & 5 \\
Londrina & 0 & 2 & 2 \\
Dourados & 1 & 2 & 3 \\
Campo Grande & 0 & 2 & 2 \\
Ambar & 0 & 2 & 2 \\
Tupy & 1 & 1 & 2 \\
Pátio de Manutenção & 2 & 0 & 2 \\
Matheus & 0 & 1 & 1 \\
Multi Aço & 0 & 1 & 1 \\
Total & 7 & 16 & 23 \\
\bottomrule
\end{tabular}
\caption{Quantidade de máquinas por local de operação.}
\end{table}

\begin{table}[H]
\centering
\small
\begin{tabular}{lrrrrr}
\toprule
Local & LIEBHERR & HYUNDAI & Total & Máq. & R\$/máq \\
\midrule
Maringá & R\$\,345.695 & R\$\,194.822 & R\$\,540.517 & 5 & R\$\,108.103 \\
Dourados & R\$\,92.821 & R\$\,304.315 & R\$\,397.137 & 3 & R\$\,132.379 \\
Prudente & R\$\,250.966 & R\$\,96.762 & R\$\,347.728 & 3 & R\$\,115.909 \\
Londrina & R\$\,0 & R\$\,200.071 & R\$\,200.071 & 2 & R\$\,100.035 \\
Ambar & R\$\,0 & R\$\,186.062 & R\$\,186.062 & 2 & R\$\,93.031 \\
Campo Grande & R\$\,0 & R\$\,169.949 & R\$\,169.949 & 2 & R\$\,84.974 \\
Pátio de Manutenção & R\$\,127.484 & R\$\,0 & R\$\,127.484 & 2 & R\$\,63.742 \\
Matheus & R\$\,0 & R\$\,84.936 & R\$\,84.936 & 1 & R\$\,84.936 \\
Tupy & R\$\,48.270 & R\$\,19.991 & R\$\,68.260 & 2 & R\$\,34.130 \\
Multi Aço & R\$\,0 & R\$\,30.636 & R\$\,30.636 & 1 & R\$\,30.636 \\
\bottomrule
\end{tabular}
\caption{Gasto por unidade de operação (soma das máquinas do local).}
\end{table}

\begin{table}[H]
\centering
\scriptsize
\begin{tabular}{lllll}
\toprule
Máquina & Marca & Local & Filial (CC) & Referência CC \\
\midrule
EHH0006 & HYUNDAI & Prudente & G\&G MATRIZ & PRESIDENTE PRUDENTE \\
EHH0009 & HYUNDAI & Prudente & G\&S PRUDENTE & PRESIDENTE PRUDENTE \\
EHL0028 & LIEBHERR & Prudente & G\&G MATRIZ & DOURADOS \\
EHH0002 & HYUNDAI & Maringá & G\&G MATRIZ & ARCELOR - PIRACICABA/SP \\
EHH0004 & HYUNDAI & Maringá & G\&G MARINGA & MARINGA DISTRITO \\
EHH0039 & HYUNDAI & Maringá & G\&G MATRIZ & PRESIDENTE PRUDENTE \\
EHL0043 & LIEBHERR & Maringá & G\&G MATRIZ & TUPY - JOINVILLE/SC \\
EHL0070 & LIEBHERR & Maringá & G\&S PRUDENTE & VEICULO, MAQUINA E EQUIPAMENTO \\
EHH0008 & HYUNDAI & Londrina & G\&G MATRIZ & LONDRINA \\
EHH0040 & HYUNDAI & Londrina & G\&G MATRIZ & LONDRINA \\
EHH0036 & HYUNDAI & Dourados & G\&G MATRIZ & DOURADOS \\
EHH0044 & HYUNDAI & Dourados & G\&G MATRIZ & TUPY - JOINVILLE/SC \\
EHL0042 & LIEBHERR & Dourados & CTA & DOURADOS \\
EHH0038 & HYUNDAI & Campo Grande & G\&G MATRIZ & CAMPO GRANDE \\
EHH0046 & HYUNDAI & Campo Grande & G\&G MATRIZ & CAMPO GRANDE \\
EHH0015 & HYUNDAI & Ambar & G\&G MATRIZ & ARCELOR - BARRA MANSA/RJ \\
EHH0019 & HYUNDAI & Ambar & G\&G MATRIZ & ARCELOR - RESENDE/RJ \\
EHH0003 & HYUNDAI & Tupy & G\&G MATRIZ & TUPY - JOINVILLE/SC \\
EHL0041 & LIEBHERR & Tupy & G\&G MATRIZ & TUPY - JOINVILLE/SC \\
EHL0013 & LIEBHERR & Pátio de Manutenção & G\&G MATRIZ & DOURADOS \\
EHL0014 & LIEBHERR & Pátio de Manutenção & G\&G MATRIZ & MARINGA DISTRITO \\
EHH0007 & HYUNDAI & Matheus & G\&G MATRIZ & PRESIDENTE PRUDENTE \\
EHH0037 & HYUNDAI & Multi Aço & G\&G MATRIZ & PRESIDENTE PRUDENTE \\
\bottomrule
\end{tabular}
\caption{Cadastro do parque: local de operação e centro de custo de referência.}
\end{table}


\section{Máquinas do parque}
\addcontentsline{toc}{section}{Máquinas do parque}

Período completo e recortes por ano (2025 e 2026 separados).

\begin{figure}[H]
  \centering
  \includegraphics[width=\linewidth]{\figdir/05_top_maquinas.png}
  \caption{Todas as máquinas do parque --- período completo (|valor\_conta|).}
\end{figure}

\begin{figure}[H]
  \centering
  \includegraphics[width=\linewidth]{\figdir/05_top_maquinas_2025.png}
  \caption{Todas as máquinas do parque --- gasto em 2025.}
\end{figure}

\begin{figure}[H]
  \centering
  \includegraphics[width=\linewidth]{\figdir/05_top_maquinas_2026.png}
  \caption{Todas as máquinas do parque --- gasto em 2026.}
\end{figure}

\subsection{Tabela completa — gasto por máquina e mês}

Anos com muitos meses usam página paisagem; períodos curtos (ex.: 2026 parcial) ficam em retrato para evitar folhas em branco.

Cada célula de mês traz o gasto total da máquina naquele período (|valor\_conta|). A coluna \textbf{Média/mês} considera somente meses com lançamento.

\begin{landscape}
\footnotesize
\setlength{\tabcolsep}{2.5pt}
\renewcommand{\arraystretch}{0.92}
\begin{longtable}{@{}ll>{\raggedright\arraybackslash}p{2.2cm}rrrrrrrrrrrrrrr@{}}
\multicolumn{18}{@{}l}{\textbf{Gasto mensal por máquina — 2025}} \\[0.2em]
\toprule
Máquina & Marca & Local & 01/25 & 02/25 & 03/25 & 04/25 & 05/25 & 06/25 & 07/25 & 08/25 & 09/25 & 10/25 & 11/25 & 12/25 & Média/mês & Total & Meses \\
\midrule
\endhead
\midrule
\multicolumn{18}{r}{\textit{Continua na próxima página}} \\
\endfoot
\bottomrule
\endlastfoot
EHH0002 & HYUNDAI & Maringá & R\$\,1,1k & R\$\,230 & --- & R\$\,1,9k & --- & --- & --- & R\$\,3,0k & R\$\,3,3k & R\$\,6,9k & --- & --- & R\$\,3,0k & R\$\,24,4k & 8 \\
EHH0003 & HYUNDAI & Tupy & R\$\,4,6k & R\$\,9,6k & R\$\,358 & R\$\,598 & --- & --- & --- & --- & --- & R\$\,1,1k & R\$\,1,9k & --- & R\$\,2,2k & R\$\,20,0k & 9 \\
EHH0004 & HYUNDAI & Maringá & R\$\,807 & --- & R\$\,2,9k & R\$\,1,4k & R\$\,7,9k & R\$\,4,0k & R\$\,1,2k & R\$\,1,9k & R\$\,13,1k & R\$\,3,0k & R\$\,32,4k & R\$\,11,6k & R\$\,6,8k & R\$\,95,0k & 14 \\
EHH0006 & HYUNDAI & Prudente & R\$\,1,6k & R\$\,3,0k & --- & --- & R\$\,166 & R\$\,2,3k & R\$\,367 & R\$\,1,2k & R\$\,2,7k & R\$\,1,9k & R\$\,350 & R\$\,179 & R\$\,1,9k & R\$\,22,4k & 12 \\
EHH0007 & HYUNDAI & Matheus & R\$\,1,6k & --- & --- & --- & R\$\,6,0k & R\$\,17,5k & R\$\,8,9k & R\$\,19,0k & R\$\,8,4k & R\$\,13,1k & R\$\,6,5k & R\$\,1,4k & R\$\,7,6k & R\$\,91,0k & 12 \\
EHH0008 & HYUNDAI & Londrina & R\$\,3,6k & R\$\,4,2k & R\$\,6,1k & R\$\,6,5k & R\$\,8,9k & R\$\,2,2k & R\$\,5,5k & R\$\,19,0k & R\$\,5,2k & R\$\,11,7k & R\$\,1,5k & R\$\,6,3k & R\$\,6,2k & R\$\,86,9k & 14 \\
EHH0009 & HYUNDAI & Prudente & --- & --- & --- & --- & R\$\,1,2k & R\$\,1,0k & R\$\,7,7k & R\$\,14,8k & R\$\,12,9k & R\$\,8,1k & R\$\,447 & R\$\,7,3k & R\$\,6,9k & R\$\,75,8k & 11 \\
EHH0015 & HYUNDAI & Ambar & R\$\,90 & --- & --- & --- & R\$\,1,9k & --- & R\$\,27,5k & R\$\,3,7k & R\$\,123 & R\$\,2,8k & R\$\,4,6k & R\$\,3,0k & R\$\,4,7k & R\$\,52,1k & 11 \\
EHH0019 & HYUNDAI & Ambar & R\$\,2,4k & R\$\,1,8k & R\$\,20,4k & R\$\,7,9k & R\$\,616 & R\$\,21,5k & R\$\,1,9k & R\$\,907 & R\$\,3,9k & R\$\,9,5k & R\$\,14,5k & R\$\,743 & R\$\,9,1k & R\$\,136,5k & 15 \\
EHH0036 & HYUNDAI & Dourados & R\$\,5,9k & R\$\,12,0k & R\$\,3,7k & R\$\,9,1k & R\$\,8,0k & R\$\,5,7k & R\$\,8,1k & R\$\,9,4k & R\$\,8,0k & R\$\,12,0k & R\$\,12,8k & R\$\,3,2k & R\$\,7,0k & R\$\,104,5k & 15 \\
EHH0037 & HYUNDAI & Multi Aço & R\$\,4,4k & R\$\,2,2k & R\$\,1,7k & R\$\,6,1k & R\$\,3,8k & R\$\,2,8k & --- & --- & --- & R\$\,278 & R\$\,1,7k & R\$\,3,0k & R\$\,2,9k & R\$\,34,4k & 12 \\
EHH0038 & HYUNDAI & Campo Grande & R\$\,13,0k & R\$\,10,0k & R\$\,20,0k & R\$\,19,7k & R\$\,7,2k & R\$\,14,0k & R\$\,16,4k & R\$\,17,9k & R\$\,18,4k & R\$\,12,7k & R\$\,8,8k & R\$\,8,9k & R\$\,12,2k & R\$\,170,5k & 14 \\
EHH0039 & HYUNDAI & Maringá & R\$\,3,7k & R\$\,641 & R\$\,3,4k & R\$\,20,8k & R\$\,4,5k & R\$\,5,9k & R\$\,4,6k & R\$\,6,5k & R\$\,15,3k & R\$\,12,0k & R\$\,4,5k & R\$\,2,4k & R\$\,5,9k & R\$\,88,7k & 15 \\
EHH0040 & HYUNDAI & Londrina & R\$\,13,0k & R\$\,5,8k & R\$\,14,7k & R\$\,6,8k & R\$\,13,2k & R\$\,9,1k & R\$\,16,1k & R\$\,10,6k & R\$\,10,0k & R\$\,15,3k & R\$\,8,3k & R\$\,2,3k & R\$\,9,0k & R\$\,135,5k & 15 \\
EHH0044 & HYUNDAI & Dourados & R\$\,11,3k & R\$\,8,0k & R\$\,12,8k & R\$\,34,9k & R\$\,20,8k & R\$\,7,1k & R\$\,1,3k & R\$\,4,1k & R\$\,78,0k & R\$\,21,0k & R\$\,13,7k & R\$\,2,7k & R\$\,15,5k & R\$\,231,9k & 15 \\
EHH0046 & HYUNDAI & Campo Grande & R\$\,2,0k & R\$\,1,6k & --- & R\$\,295 & R\$\,11,9k & --- & --- & --- & R\$\,1,2k & R\$\,2,6k & --- & --- & R\$\,2,9k & R\$\,20,3k & 7 \\
EHL0013 & LIEBHERR & Pátio de Manutenção & --- & R\$\,149 & --- & --- & --- & R\$\,1,8k & --- & R\$\,713 & R\$\,8,7k & --- & R\$\,416 & --- & R\$\,2,1k & R\$\,14,4k & 7 \\
EHL0014 & LIEBHERR & Pátio de Manutenção & --- & R\$\,725 & --- & --- & --- & R\$\,11,8k & R\$\,21,8k & R\$\,15,9k & R\$\,18,0k & R\$\,6,6k & R\$\,3,8k & R\$\,2,1k & R\$\,10,3k & R\$\,113,1k & 11 \\
EHL0028 & LIEBHERR & Prudente & R\$\,6,9k & R\$\,6,1k & R\$\,14,9k & R\$\,14,6k & R\$\,15,6k & R\$\,11,4k & R\$\,12,7k & R\$\,58,8k & R\$\,22,9k & R\$\,34,8k & R\$\,14,0k & R\$\,11,1k & R\$\,16,8k & R\$\,251,3k & 15 \\
EHL0041 & LIEBHERR & Tupy & R\$\,7,8k & R\$\,1,6k & R\$\,215 & R\$\,8,2k & --- & R\$\,7,6k & R\$\,1,5k & R\$\,660 & R\$\,3,1k & R\$\,6,2k & --- & R\$\,3,1k & R\$\,3,7k & R\$\,48,3k & 13 \\
EHL0042 & LIEBHERR & Dourados & --- & --- & --- & --- & --- & R\$\,1,7k & R\$\,5,6k & R\$\,6,5k & R\$\,5,0k & R\$\,9,5k & R\$\,32,9k & R\$\,14,0k & R\$\,9,3k & R\$\,92,8k & 10 \\
EHL0043 & LIEBHERR & Maringá & R\$\,6,1k & R\$\,9,6k & R\$\,246 & R\$\,5,0k & R\$\,2,9k & R\$\,14,4k & R\$\,11,4k & R\$\,188,0k & R\$\,15,4k & R\$\,18,4k & R\$\,26,4k & R\$\,4,2k & R\$\,23,2k & R\$\,348,2k & 15 \\
EHL0070 & LIEBHERR & Maringá & --- & --- & --- & --- & --- & --- & --- & --- & --- & --- & R\$\,400 & --- & R\$\,400 & R\$\,400 & 1 \\
\end{longtable}
\end{landscape}

\begin{table}[H]
\centering
\footnotesize
\setlength{\tabcolsep}{3.5pt}
\renewcommand{\arraystretch}{0.95}
\caption{Gasto mensal por máquina — 2026}
\begin{tabular}{@{}ll>{\raggedright\arraybackslash}p{2.2cm}rrrrrr@{}}
\toprule
Máquina & Marca & Local & 01/26 & 02/26 & 03/26 & Média/mês & Total & Meses \\
\midrule
EHH0002 & HYUNDAI & Maringá & --- & R\$\,6,1k & R\$\,1,8k & R\$\,3,0k & R\$\,24,4k & 8 \\
EHH0003 & HYUNDAI & Tupy & R\$\,155 & R\$\,1,3k & R\$\,463 & R\$\,2,2k & R\$\,20,0k & 9 \\
EHH0004 & HYUNDAI & Maringá & R\$\,5,8k & R\$\,3,1k & R\$\,6,0k & R\$\,6,8k & R\$\,95,0k & 14 \\
EHH0006 & HYUNDAI & Prudente & R\$\,108 & --- & R\$\,8,5k & R\$\,1,9k & R\$\,22,4k & 12 \\
EHH0007 & HYUNDAI & Matheus & R\$\,102 & R\$\,7,0k & R\$\,1,6k & R\$\,7,6k & R\$\,91,0k & 12 \\
EHH0008 & HYUNDAI & Londrina & --- & R\$\,4,7k & R\$\,1,6k & R\$\,6,2k & R\$\,86,9k & 14 \\
EHH0009 & HYUNDAI & Prudente & R\$\,1,1k & R\$\,8,6k & R\$\,12,7k & R\$\,6,9k & R\$\,75,8k & 11 \\
EHH0015 & HYUNDAI & Ambar & R\$\,1,0k & R\$\,2,0k & R\$\,5,3k & R\$\,4,7k & R\$\,52,1k & 11 \\
EHH0019 & HYUNDAI & Ambar & R\$\,39,1k & R\$\,8,8k & R\$\,2,3k & R\$\,9,1k & R\$\,136,5k & 15 \\
EHH0036 & HYUNDAI & Dourados & R\$\,581 & R\$\,1,3k & R\$\,4,7k & R\$\,7,0k & R\$\,104,5k & 15 \\
EHH0037 & HYUNDAI & Multi Aço & R\$\,1,7k & R\$\,1,5k & R\$\,5,3k & R\$\,2,9k & R\$\,34,4k & 12 \\
EHH0038 & HYUNDAI & Campo Grande & R\$\,2,9k & --- & R\$\,719 & R\$\,12,2k & R\$\,170,5k & 14 \\
EHH0039 & HYUNDAI & Maringá & R\$\,2,3k & R\$\,1,1k & R\$\,1,2k & R\$\,5,9k & R\$\,88,7k & 15 \\
EHH0040 & HYUNDAI & Londrina & R\$\,613 & R\$\,4,3k & R\$\,5,3k & R\$\,9,0k & R\$\,135,5k & 15 \\
EHH0044 & HYUNDAI & Dourados & R\$\,836 & R\$\,8,6k & R\$\,6,5k & R\$\,15,5k & R\$\,231,9k & 15 \\
EHH0046 & HYUNDAI & Campo Grande & --- & R\$\,675 & --- & R\$\,2,9k & R\$\,20,3k & 7 \\
EHL0013 & LIEBHERR & Pátio de Manutenção & R\$\,1,4k & --- & R\$\,1,3k & R\$\,2,1k & R\$\,14,4k & 7 \\
EHL0014 & LIEBHERR & Pátio de Manutenção & R\$\,15,3k & R\$\,4,8k & R\$\,12,2k & R\$\,10,3k & R\$\,113,1k & 11 \\
EHL0028 & LIEBHERR & Prudente & R\$\,5,3k & R\$\,15,0k & R\$\,7,3k & R\$\,16,8k & R\$\,251,3k & 15 \\
EHL0041 & LIEBHERR & Tupy & R\$\,5,1k & R\$\,2,5k & R\$\,669 & R\$\,3,7k & R\$\,48,3k & 13 \\
EHL0042 & LIEBHERR & Dourados & R\$\,2,1k & R\$\,2,1k & R\$\,13,6k & R\$\,9,3k & R\$\,92,8k & 10 \\
EHL0043 & LIEBHERR & Maringá & R\$\,18,3k & R\$\,3,9k & R\$\,24,0k & R\$\,23,2k & R\$\,348,2k & 15 \\
EHL0070 & LIEBHERR & Maringá & --- & --- & --- & R\$\,400 & R\$\,400 & 1 \\
\bottomrule
\end{tabular}
\end{table}



\subsection{Gasto mensal por máquina (top 5)}

\begin{figure}[H]
  \centering
  \includegraphics[width=\linewidth]{\figdir/09_mes_maquinas_liebherr.png}
  \caption{Liebherr --- gasto mensal das 5 máquinas com maior volume.}
\end{figure}

\begin{figure}[H]
  \centering
  \includegraphics[width=\linewidth]{\figdir/09_mes_maquinas_hyundai.png}
  \caption{Hyundai --- gasto mensal das 5 máquinas com maior volume.}
\end{figure}

\section*{Conclusão}
\addcontentsline{toc}{section}{Conclusão}

Em síntese, as duas marcas apresentam perfis distintos de custo. O comparativo principal usa o custo médio por máquina do parque (volume total $\div$ 7 ou 16), não o volume bruto --- Liebherr tem 7 máquinas e Hyundai 16, e totais absolutos enviesam a leitura. Diferem na intensidade por máquina e no mix entre manutenção, combustível e peças. Maringá — R\$ 540,517 (5 máq.; R\$ 108,103/máq). Recomenda-se priorizar revisão dos ativos líderes de gasto e acompanhar a evolução mensal para validar ações de contenção de custo.


\vfill
{\small\textit{Fonte:} \texttt{base\_maquinas\_2025\_2026}. Parque: 7 Liebherr, 16 Hyundai.
Comparativos entre marcas em custo médio por máquina do parque.
Emitido em \reportdate.}

\end{document}
